% Proposed section for the d-SEAMS 2.0 paper.
% Numbers are from single trajectories. The slope error is the
% standard error of an ordinary-least-squares fit along that trajectory.

\section{Salt at a Madrid-2019 freezing front}
\label{sec:brine-science}

The electrolyte papers of the Madrid group measure a coexistence
molality, not the structure of the moving front. Lamas, Vega and
Noya put ice Ih in contact with a Madrid-2019 solution and followed
the liquid molality until it stopped changing.\cite{lamas2022}
For NaCl at $1\,\mathrm{bar}$ their equilibrium molalities are
$1.19(0.05)\,\mathrm{mol\,kg^{-1}}$ at $247\,\mathrm{K}$,
$2.13(0.08)\,\mathrm{mol\,kg^{-1}}$ at $243\,\mathrm{K}$,
$4.07(0.08)\,\mathrm{mol\,kg^{-1}}$ at $236\,\mathrm{K}$ and
$6.24(0.05)\,\mathrm{mol\,kg^{-1}}$ at $224\,\mathrm{K}$, from runs
of $1.1$--$1.8\,\mu\mathrm{s}$.
Blazquez, Conde and Vega repeated the $236\,\mathrm{K}$ point for
several scaled charges and obtained
$4.17\,\mathrm{mol\,kg^{-1}}$ for Madrid-2019
($q=\pm 0.85$) against an experimental $3.68\,\mathrm{mol\,kg^{-1}}$;
the $q=\pm 0.75$ Madrid-Transport model gave $3.73$.\cite{blazquez2023}
Soria, Espinosa, Ramirez, Valeriani, Vega and Sanz measured
nucleation rates, not a planar front, and they used Joung--Cheatham
ions with TIP4P/2005 rather than Madrid-2019.\cite{soria2018}
At $1.85\,\mathrm{mol\,kg^{-1}}$ their direct-coexistence melting
point is $239\pm 1\,\mathrm{K}$, and they take $250\,\mathrm{K}$ as
the melting point of pure TIP4P/2005.
Goswami, Blazquez, Fern{\'a}ndez-Sedano, Noya, J{\'o}nsson, Troncoso
and Vega extended Madrid-2019 to $\mathrm{Fe^{2+}}$ and used the
viscosity to read the extent of complexation in $\mathrm{FeCl_2}$
and $\mathrm{MgCl_2}$.\cite{goswami2026}
That paper does not freeze the solution.
No paper coauthored by Amrita Goswami and Eduardo Sanz appears in
the arXiv or OpenAlex lists consulted for this section.
Conde, Rovere and Gallo did grow a TIP4P/2005 front from a salty
liquid, on the microsecond scale, and found that chloride replaces
two water molecules while sodium sits in an interstitial
site.\cite{conde2017}
The ion model in that work is an earlier unit-charge parameterisation
whose solubility they quote as $5.8\,\mathrm{mol\,kg^{-1}}$, not the
scaled Madrid-2019 charges.
None of these studies reports a CHILL+ rejection coefficient, the
ring census of the front layer, or the size of a brine pocket.

The deck used here is TIP4P/2005 with the Madrid-2019 parameters of
Zeron, Abascal and Vega, charges $q=\pm 0.85$, and the tabulated
cross terms rather than Lorentz--Berthelot mixing.\cite{zeron2019}
Madrid-Transport was not substituted for it: the freezing-point
papers above, and the $\mathrm{FeCl_2}$ extension, are Madrid-2019.
As a check on the ion parameters, a $4\times 4\times 4$ rock-salt
crystal ($512$ ions) was run in the $NpT$ ensemble at $298.15\,\mathrm{K}$
and $1\,\mathrm{bar}$ for $8\,\mathrm{ps}$ after a minimisation.
The density over the second half is
$2.044\pm 0.001\,\mathrm{g\,cm^{-3}}$
(six blocks).
Table~V of Zeron, Abascal and Vega gives
$2.050\,\mathrm{g\,cm^{-3}}$ for the same scaled-charge crystal.
The residual is $0.3\,\%$. Their average was $10\,\mathrm{ns}$ on
$1000$ ions in GROMACS, so the comparison is only at that level.

The growth cell is $128$ ice Ih molecules and $128$ liquid waters
($2\times 2\times 2$ orthogonal cells; cross section
$15.7\times 14.8\,\mathrm{\AA}$).
Hydrogens on the ice satisfy the ice rules; the liquid and the ions
occupy the opened $z$ gap.
The nominal molality is that of the liquid half.
Each trajectory is one run of $0.50\,\mathrm{ns}$ after
$10\,\mathrm{ps}$ with the ice pinned, in the $NpT$ ensemble with
only $z$ barostatted, at $1\,\mathrm{bar}$.
Frames are every $10\,\mathrm{ps}$ ($51$ frames) and are not stored
in the repository; \texttt{input/brine/grow.py} regenerates them.
Ice is the published CHILL+ label on the oxygens (hexagonal or cubic).
The front position is the $z$ at which the occupied-bin ice fraction
crosses one half. Empty bins, which are the gaps between ice planes,
do not split the slab.

\begin{table}[t]
\centering
\caption{Single-trajectory CHILL+ front on the Madrid-2019 slab.
$v$ is the slope of the ice thickness per front
(half the slope of the slab thickness).
$\mathrm{d}N/\mathrm{d}t$ is the slope of the number of crystal
oxygens. Both uncertainties are the standard error of the slope
along that trajectory. $c_L$ is the liquid molality in the last
frame. Defects and ring counts are means over the front layer.
No ion entered a bulk-ice bin, so $k=c_\mathrm{ice}/c_L=0$, and
there were no trapping transitions.}
\label{tab:brine}
\begin{tabular}{lrrrrrr}
\hline
$m_0$ & $T$ & $v$ & $\mathrm{d}N/\mathrm{d}t$ & $c_L$ & defects & 6-rings \\
$\mathrm{mol\,kg^{-1}}$ & K & \AA\,ns$^{-1}$ & ns$^{-1}$ & $\mathrm{mol\,kg^{-1}}$ & & \\
\hline
0   & 220 & $1.39\pm 0.39$ & $28\pm 4$  & 0    & 7.8 & 72 \\
0   & 236 & $1.10\pm 0.53$ & $15\pm 6$  & 0    & 7.8 & 75 \\
0   & 240 & $4.74\pm 0.56$ & $61\pm 6$  & 0    & 9.1 & 75 \\
0   & 255 & $3.62\pm 0.59$ & $41\pm 6$  & 0    & 11.5 & 76 \\
0.5 & 220 & $-0.19\pm 0.38$ & $2\pm 3$ & 0.43 & 9.1 & 74 \\
1   & 220 & $-3.10\pm 0.49$ & $-40\pm 5$ & 0.63 & 11.7 & 78 \\
1   & 236 & $-2.71\pm 0.57$ & $-40\pm 6$ & 0.72 & 11.0 & 76 \\
2   & 220 & $-1.14\pm 0.44$ & $-11\pm 4$ & 1.77 & 11.2 & 78 \\
\hline
\end{tabular}
\end{table}

On this slab the pure-water front advances.
One {\AA}\,ns$^{-1}$ is $10\,\mathrm{cm\,s^{-1}}$, so the
$220\,\mathrm{K}$ pure-water front moves at
$14\pm 4\,\mathrm{cm\,s^{-1}}$ and the crystal count rises by
$28\pm 4$ oxygens per nanosecond.
Adding salt at the same temperature stops that advance.
At $0.5\,\mathrm{mol\,kg^{-1}}$ the crystal-count slope is
$2\pm 3\,\mathrm{ns^{-1}}$, indistinguishable from zero on this
trajectory.
At $1$ and $2\,\mathrm{mol\,kg^{-1}}$ the CHILL+ crystal shrinks
($-40\pm 5$ and $-11\pm 4$ oxygens per nanosecond).
The ordering of those two salted slopes is a property of these two
trajectories, which contain four and ten ions.
It is not a replicate estimate.
The same shrinkage appears at $1\,\mathrm{mol\,kg^{-1}}$ and
$236\,\mathrm{K}$.
Lamas, Vega and Noya find that a $1$--$2\,\mathrm{mol\,kg^{-1}}$
solution is still below the equilibrium molality at
$224$--$236\,\mathrm{K}$, so a microsecond run should freeze, not
melt.\cite{lamas2022}
The loss of CHILL+ crystal seen here in $0.5\,\mathrm{ns}$ is the
disorder of the front, not that coexistence molality.
The pure-water run at $255\,\mathrm{K}$, five kelvin above the
$250\,\mathrm{K}$ melting point used by that paper, still gains
crystal oxygens.
Half a nanosecond on a $15\,\mathrm{\AA}$ cross section does not
locate the melting point.

Every salted frame has $k=0$: no ion sits in a bulk-ice bin, and
the ion state never changes from liquid to ice or to the front
shell.
With zero trapping events in $0.5\,\mathrm{ns}$, incorporation is
rarer than the microsecond chloride substitution reported for the
unit-charge model.\cite{conde2017}
The liquid molality at the end of the salted runs lies below the
prepared value when the crystal count falls, because water has
joined the liquid.
Five-membered rings in the front layer stay near six per frame
while six-membered rings stay near $75$.
The mean number of coordination defects in that layer rises from
$7.8$ in pure water at $220\,\mathrm{K}$ to $11.7$ and $11.2$ at
$1$ and $2\,\mathrm{mol\,kg^{-1}}$.
Brine pockets, defined as liquid-water components that are not the
reservoir in contact with the front, are rare.
The mean count per frame is between $0.04$ and $0.24$, and the
pockets that do appear contain one or two waters.
A population of nanometre brine pockets is not present in this cell.

On the first frame of the $220\,\mathrm{K}$ pure-water run, freud's
Steinhardt $q_6$ (cutoff $3.5\,\mathrm{\AA}$) has median $0.628$ on
the $98$ oxygens CHILL+ calls crystal and median $0.526$ on the
others.
The crystal values sit on a single peak: the $10$th and $90$th
percentiles are $0.6279$ and $0.6285$.
A pre-set threshold $q_6>0.45$ labels every crystal oxygen and also
$118$ non-crystal oxygens, so the two assignments agree on $138$ of
$256$ oxygens.
freud does not implement CHILL+; the check is that the crystal
label is the sharp high-$q_6$ peak, while a threshold below the
liquid median is not.

The analysis time is for one frame of this trajectory, serial, with
the neighbour list, published CHILL+ and the front measurements
together.
Eight repeats give
$9.30$, $9.00$, $8.40$, $8.51$, $8.54$, $8.69$, $8.97$ and
$8.58\,\mathrm{ms}$.
The median is $8.63\,\mathrm{ms}$.

\begin{thebibliography}{9}
\bibitem{lamas2022}
C.~P. Lamas, C.~Vega and E.~G. Noya,
J.~Chem.\ Phys.\ \textbf{156}, 134503 (2022),
doi:10.1063/5.0085051.
\bibitem{blazquez2023}
S.~Blazquez, M.~M. Conde and C.~Vega,
J.~Chem.\ Phys.\ \textbf{158}, 054505 (2023),
doi:10.1063/5.0136498.
\bibitem{soria2018}
G.~D. Soria, J.~R. Espinosa, J.~Ramirez, C.~Valeriani, C.~Vega
and E.~Sanz,
J.~Chem.\ Phys.\ \textbf{148}, 222811 (2018),
doi:10.1063/1.5008889.
\bibitem{goswami2026}
A.~Goswami, S.~Blazquez, L.~Fern{\'a}ndez-Sedano, E.~G. Noya,
H.~J{\'o}nsson, J.~Troncoso and C.~Vega,
J.~Phys.\ Chem.~B (2026),
doi:10.1021/acs.jpcb.6c01329.
\bibitem{conde2017}
M.~M. Conde, M.~Rovere and P.~Gallo,
Phys.\ Chem.\ Chem.\ Phys.\ \textbf{19}, 9566 (2017),
doi:10.1039/c7cp00665a.
\bibitem{zeron2019}
I.~M. Zeron, J.~L.~F. Abascal and C.~Vega,
J.~Chem.\ Phys.\ \textbf{151}, 134504 (2019),
doi:10.1063/1.5121392.
\end{thebibliography}
